% sections/discussion.tex -- Discussion (W25): executor-view contract, IETF htm/htu, limitations, ethics.
\section{Discussion}\label{sec:discussion}

\paragraph{An executor-view contract.}\label{sec:discussion:rec} RFC~9449 tells
the resource server to compare \htm and \htu with ``the HTTP method'' and ``the
HTTP URI value'' of the request~\cite{rfc9449}. That is unambiguous for one
process and underspecified for a layered deployment. Our results suggest
stating \emph{whose} request is meant: the method and path the
\emph{executor} acts on, after override resolution and routing
normalisation. A verifier placed elsewhere must either receive exactly the
representation the gateway forwards upstream, or hand the executor a decision
bound to its view that the executor checks against what it routed (D-A, D-B).
Re-verifying the same raw view at the origin is not enough (\Abl{3}). A
rewriting gateway is a party to the contract: it must publish its
public-to-internal path map, or the verifier must bind the path the gateway
forwards; otherwise the \FRRAllDAB{} rewrite false rejections of
\S\ref{sec:defence} are the price of soundness.

\paragraph{The \htm/\htu proposal.} Proposals to make \htm and \htu optional
because proxies rewrite URIs~\cite{allam2026thread,madden2025thread} would turn
every \GB witness into a \GA one: in \Abl{1}, dropping \htu more than doubles
the hits and admits \AblOneSwapNoHtu{} cross-operation replays, and dropping
\htm admits \AblOneSwapNoHtm{} method swaps that a mutation corpus cannot see.
Rewriting is a reason to define the executor's view, not to drop the binding.
Where operations dispatch on the query or body (action endpoints, RPC, FHIR
conditional interactions, MCP tool calls), an optional query claim and a
digest checked by a body-holding hop close \GA dispatch; header-carried
semantics and \GB remain, so binding more is complementary to the contract,
not a substitute.

\paragraph{Gateways and frameworks.} Every gateway-resolved class has one
shape: the authorization subrequest carries one representation of the path and
the upstream request another. HAProxy, which uses one parsed URI for both, had
no \GB hit. Gateways should authorize what they forward, reject a \texttt{\#}
in an origin-form target as RFC~9112 requires, and strip client-supplied
\hdr{X-Verified-*} and certificate headers. Method override should be off by
default and visible to authorization middleware; Laravel offers no switch, and
Rails keeps it in every non-API application.

\paragraph{Limitations.}\label{sec:discussion:lim}
\emph{Route space.} The bench API has \NumBenchOps{} operations plus
\NumNewRoutes{} in v2; several differentials change the operation only when a
matching route exists (nested route for B1, shared POST path for override,
catch-all for B3). Whether real applications contain such routes is argued
from \Exp{1}/\Exp{2}, not measured. The S1 rule also counts identifier
suffixes that reach no other resource.
\emph{APISIX mechanism.} The verifier-side half of B2 is unconfirmed: the
published \texttt{forward-auth} code should send the raw URI, which conflicts
with the decoded view we logged (\S\ref{sec:results}).
\emph{Testbed.} Every topology ran on one 16\,GB laptop over cleartext
HTTP/1.1 with versions pinned on one date (Kong 3.9.5 is not its newest
line; the split is configuration-level); TLS, HTTP/2 front ends and
multi-replica replay caches were not exercised, and only depth-1 cases ran.
\emph{False rejection} was measured behind NGINX only; other gateways' FRR is
unmeasured. \emph{Latency} is from the shared laptop only (a dedicated-server
re-run did not complete), so it supports ``negligible at \PerfRate{}\,req/s''
and nothing about throughput.
\emph{Coverage.} \GC appears only as propagation because the bench API does
not act on trust headers; the \Aclient battery omitted absolute-form targets
and \hdr{X-Original-URL}, and ran \NumEFourTopologies{} of \NumEFourCells{}
cells. D-B was implemented for four of eight frameworks. One author screened
the literature and judged the prevalence sample.

\paragraph{Ethics and disclosure.}\label{sec:discussion:ethics} All experiments
ran against software we deployed ourselves, on local containers with public
release images, test keys and synthetic accounts; no third-party or production
system was tested. We drafted reports for \NumDisclosures{} candidates and
sent none. The \texttt{\#} handling of uvicorn/h11, Go \texttt{net/http} and
FrankenPHP is known but unfixed, and will be reported to the projects before
publication. After the prior-art check we no longer treat the APISIX
observation as new: if reported, it will be framed as a consequence of the
intentional CVE-2026-124680 behaviour, after the mechanism is confirmed.
Configuration-level and by-design behaviours (Kong raw variables, default-off
normalisation, method override) are documented and appear as recommendations.
